% !PRINTED none
\weberpage{none}
\begin{center}
\textbf{\gesperrt{Berichtigungen zum ersten Bande.}}
\end{center}

\begin{center}
\rule{4cm}{0.4pt}
\end{center}

\begin{flushleft}
Seite 93, Zeile 6 v.\,u. lies $B_h^{(i)}$ statt $B^{(i)}$. \\
„ 117, Formel (10) lies $\sqrt[3]{\dfrac{-b}{2}+\sqrt{R}}$ statt $\sqrt{\dfrac{-b}{2}+\sqrt{R}}$. \\
„ 158, Formeln (14) rechts $m$ und $n$ zu vertauschen. \\
„ 172, Formel (14) lies $n-1$ statt $n+1$. \\
„ 185, Zeile 9 v.\,o. „ $\Phi'$ statt $\Phi$. \\
„ 190, „ 11 v.\,u. „ $a_1'$ „ $(n-1)a_1'$. \\
„ 194, „ 7 v.\,u. „ $+9$ „ $-9$. \\
„ 200, „ 7 v.\,u. „ \S.\,61 „ \S.\,67. \\
„ 202, „ 3 v.\,u. „ $u$ „ $v_1$ (in beiden Formeln). \\
„ 213, „ 6 v.\,u. „ $\Phi(\tau,\xi)$ statt $\Phi(\tau,k)$. \\
„ 221, „ 12 v.\,u. „ $\dfrac{2}{9}$ statt $\dfrac{2}{3}$. \\
„ 222, „ 4,5 v.\,u. „ $+q_0, -q_0$ statt $-q_0, +q_0$. \\
„ 233, „ 4,7 v.\,o. „ $\tfrac{1}{5}D$ statt $5D$. \\
„ 236, „ 12 v.\,o. „ $5c^2$ „ $-5c^2$. \\
„ 251, „ 10 v.\,o. „ $-\psi_3$ „ $\psi_3$. \\
„ 268, „ 5 v.\,u. „ $\tfrac{1}{5}D$ „ $D$. \\
„ 282, Formel (2) „ $f_s$ „ $f^s$. \\
„ 318, Zeile 2 v.\,u. „ $F(z)$ statt $F(x)$. \\
„ 332, Formel (5) „ $f(\beta)-$ statt $f(\beta)+$. \\
„ 334, Zeile 6 v.\,u. „ 0,00164 „ 0,0164. \\
„ 344, „ 6 v.\,u., Formel, lies $\alpha_n\varphi(\alpha_n)$ und $\varphi(\alpha_n)$ statt $\alpha_m\varphi(\alpha_m)$
in $\varphi(\alpha_m)$. \\
„ 357, „ 3 v.\,u. lies 6 statt 5. \\
„ 378, 379, Formel (13), (14) ist $a_1, b_1, c_1$ durch $\pm a_1, \pm b_1, \pm c_1$ zu
ersetzen, und das Zeichen so zu bestimmen, dass $c_1$ positiv wird. \\
„ 402, Zeile 6,7 v.\,u. ist $a_4=1$ und $a_3=2$ zu vertauschen. \\
„ — „ 4 v.\,u. lies 9083 statt 9183. \\
„ 419, Hierzu vergleiche man den I. Nachtrag zum zweiten Bande. \\
„ 423, Zeile 15 v.\,o. lies $\dfrac{(n-1)(n-2)}{2}$ statt $\dfrac{n\cdot n-1}{2}$. \\
„ — „ 24 v.\,o. ist zu lesen
$\dfrac{1}{2}\left(\dfrac{(n-1)(n-2)}{2}-\dfrac{n-1}{2}\right)
=\dfrac{(n-1)(n-3)}{4}\equiv 0\ (\mathrm{mod}\ 2)$.
\end{flushleft}
